\section{结论}
\label{sec:conclusion}

本文围绕函数级代码生成中的测试时计算组织问题，提出预算分档编排协议。该协议在统一执行器、统一判定口径和统一令牌记账下，将流程固定为单步、单步+修补、分流+重链与救援、级联四个可比较档位，把算得更多转化为预算投向何处的可复现实验问题。

HumanEval 与 MBPP 双基准上的三随机种子主结果显示，分档编排形成了清晰的成本--收益前沿：MBPP 通过率从 $82.51\%$ 提升至 $91.72\%$，HumanEval 从 $97.56\%$ 提升至 $99.80\%$。机制实验表明，结构难度分适合作为预算先验，压缩反馈是经济有效的回注形式，级联档的主要收益来自候选扩展。

总体而言，预算分档编排提供了一种面向部署的代码生成研究口径。中预算优先修补与结构化主链，高预算优先候选扩展，是当前数据下最有解释力的策略。后续工作将扩展多随机种子路由对照、跨模型版本回归以及更贴近真实部署成本的联合记账模型。
